\section{Weekly Summary}

\begin{tcolorbox}[colback=yellow!10,colframe=black!50,title=AI-Generated Content]
\noindent The summary below was written by Claude (Anthropic) from that week's regression outputs and series descriptions. It has not been reviewed or fact-checked by a human, and should be read as an automated, informal recap -- not analysis.
\end{tcolorbox}

Welcome back to the weekly parade of statistical nonsense, where we pair up economic time series like a blindfolded matchmaker and then act shocked when they don't get along. This week featured a truly inspired roster of contestants, roughly half of which carry the word "DISCONTINUED" in their titles, which really sets the tone. Nothing says rigorous science like regressing a dead series against another dead series and nodding thoughtfully.

Monday kicked things off with Mid-Atlantic residential natural gas prices versus the percent of syndicated loans at foreign bank branches, because why not. The raw regression came back with an R-squared of 0.44 and a p-value of 0.001, which in a less self-aware newsletter would be the point where we announce a breakthrough. Alas, the moment we detrended, the whole thing dissolved into a puddle (R-squared 0.02, p-value 0.56, coefficient sheepishly backing away from the room). Classic trend mirage. Two things drifted over the same five years, held hands briefly, and were revealed to be total strangers the instant we made them stand on their own. Tuesday's matchup of bank equity ratios against the dollar-to-SDR exchange rate pulled the same stunt: "significant" raw (p=0.003) and then a limp, insignificant detrended shrug (p=0.46). Heartbreaking. Not really.

Then came the genuinely interesting ones, and I cannot believe I'm typing that. Wednesday paired Hong Kong's holdings of U.S. long-term corporate bonds with, and I want to stress this is real, commercial carbon dioxide emissions from natural gas pipelines. The raw regression was significant, sure, but so was the detrended one (p=0.003), and the relationship even flipped sign once we stripped the trends out. That is the rare case where the correlation survives having its training wheels removed. Do I have a theory? Absolutely: bonds and pipeline emissions are clearly the same thing viewed from different dimensions. Do I believe that theory? No. Friday's furniture-and-appliance retail inventories versus New England recreational goods spending did the same honorable thing, staying significant in both raw (p=0.004) and detrended (p=0.01) forms, which at least makes vague economic sense since both involve people buying durable stuff. Mark your calendars, two survivors in one week.

The rest of the week returned us to our regularly scheduled disappointment: Thursday's SBA-backed loan percentages versus more bank equity ratios never showed up at all (p=0.17 raw, p=0.20 detrended, a committed non-relationship), and Saturday's unemployment-versus-employment matchup posted a magnificent R-squared of 0.0002, which is statistics for "these two have never met." As always, the standard disclaimer applies with extreme vigor: these are randomly paired series, I am drawing exactly zero real conclusions, and spurious correlation is not the exception here but the entire business model. The fact that two pairs held up under detrending this week is fun, not meaningful, and if anyone cites the bond-pipeline finding in a serious context, that's between them and their conscience.
